\section{The actual cubic diagonal point is unique}
\label{sec:cubic-fixed-point}

Keep $K=\mathbb Q(\omega)$, the actual cubic seed and
normalized ring of AC-G49, and the coefficient-prime
range of that result.  In particular $\ell\ne2,3,269$,
the old \texttt{GD} hypotheses hold, and the entire
selected local functor is transported by the constant
diagonal matrix of AC-G49-inert.
The assertions here use that transport, including
the ordered ordinary inertia conclusion for all its
coefficient-field points.  They do not claim this
conclusion for a newly constructed local prime merely
from a dimension or Hodge-type calculation.

\begin{theorem}[AC-G55: the fixed-point field for the actual cubic ring]
\label{AC-G55}
Let $f:S\to\overline{\mathbb Q}_\ell$ be an actual
\texttt{IsLocPt} of the ring of AC-G49.  If
\[
 \texttt{Upper}(\eta,\operatorname{push}(f,r^{\rm univ}))
 \quad\hbox{and}\quad
 \texttt{Lower}(\eta,\operatorname{push}(f,r^{\rm univ})),
                                                               \tag{G55.1}
\]
then
\[
              \operatorname{push}(f,r^{\rm univ})=r_\chi,
                \qquad f=\overline{x}_0.                      \tag{G55.2}
\]
Thus the \texttt{PtHyp.fixed} field in
\path{UnitaryThreeGroups/Rigid0.lean} holds for this
actual cubic ring.  No automorphy theorem or
nonbase-point tangent vanishing is required.
\end{theorem}
\begin{proof}
Put $r=\operatorname{push}(f,r^{\rm univ})$.
By the definitions in
\path{UnitaryThreeGroups/BaseRepGL3NormalisedOfPsiEqOne.lean},
(G55.1) says that all positive- and negative-weight
entries are zero on $H=\ker\eta=G_K$.
Consequently
\[
                  r(h).{\rm mat}
                  =\operatorname{diag}(a_0(h),a_1(h),a_2(h))
                              \qquad(h\in H),                 \tag{G55.3}
\]
where the $a_i$ are continuous characters.
The multiplier on $H$ is $1$, so these are ordinary
multiplicative characters.
The actual normalization is $r(c).{\rm mat}=I$.
Every $g\in\Gamma$ is $h$ or $hc$, and the semidirect
product law therefore makes its matrix diagonal too.

The finite-coefficient and integrality statements
are valid for this actual point.
Apply \texttt{Loc.ptFin} in
\path{PadicHodgeTheory/PtFinFinite.lean} to its constant
dual-number lift; the entire image lies in one finite
extension $L/\mathbb Q_\ell$.
For each $y\in S$, choose $b\in O$ having the same
residue.  Then $y-b\in\mathfrak m_S$, so
$|f(y)-b|_\ell<1$ and $|f(y)|_\ell\le1$.
Diagonal entries and their inverses are images of
elements of $S$, and are consequently units of
$\mathcal O_L$.  Their reductions equal the
prescribed residual diagonal entries.
In particular the three characters
\[
 b_0=\chi^j\gamma|_H,\qquad b_1=1,\qquad
 b_2=\chi^j\gamma^{-1}|_H,\qquad
                     \delta_i=a_i b_i^{-1}                   \tag{G55.4}
\]
take values in $\mathcal O_L^\times$, and each
$\delta_i$ reduces to $1$.

\proofstep{Every prime away from the coefficient prime.}
Let $w\nmid\ell$ be a finite place of $K$, and choose
its inertia subgroup $I_w\subset H$, inside an inertia
subgroup of $\Gamma$ above the underlying rational
prime.  Put $m=3\operatorname{ord}(\psi)$.
The exact \texttt{GD.ord} field in
\path{HeightAndVanishing/Counts.lean} says
$\ell\nmid\operatorname{ord}(\psi)$; hence $\ell\nmid m$.
The cyclotomic character is trivial on $I_w$, while
$\chi^3=1$ and $\psi^{\operatorname{ord}(\psi)}=1$.
Thus
\[
              b_i(g)^m=1,\qquad
              \overline r_\chi(g)^m=1
                           \qquad(g\in I_w).                 \tag{G55.5}
\]
For any such $g$, its power $g^m$ belongs to the
actual inertia kernel killed by the ring.
Therefore $r(g)^m=r(g^m)=1$, giving
$a_i(g)^m=1$ and $\delta_i(g)^m=1$.
The reduction of $\delta_i(g)$ is $1$.
Since
\[
 X^m-1=(X-1)(1+X+\cdots+X^{m-1})
\]
and the second factor at $\delta_i(g)$ reduces to
$m\ne0$ in the residue field, that second factor is
a unit.  It follows that $\delta_i(g)=1$.
This proves unramifiedness at every $w\nmid\ell$,
including the primes above $3$ and $269$.
There is no assumption that $\psi$ is unramified at
$269$.

\proofstep{The coefficient-prime inertia and its ordering.}
On coefficient-prime inertia, $\eta=1$ and $\chi=1$.
The constant change of local frame in AC-G49-inert is
diagonal; it commutes with (G55.3) on inertia.
After that change of frame the whole local point
belongs to the old selected prime.
Apply
\path{PadicHodgeTheory/CommutingUpperTriangularDistinct.lean},
\texttt{Loc.lemma2b}.
Its conclusion is an actual ordered inertia flag
with characters
\[
                         \gamma,\quad1,\quad\gamma^{-1}.
                                                               \tag{G55.6}
\]
This uses the full selected-prime point, not a
semisimplified Weil--Deligne type alone.

The exact \texttt{dat\_facts} in
\path{UnitaryThreeGroups/DatFacts.lean} gives
$s$ in that inertia group with
$\gamma(s)^2-1\in O^\times$.
Thus the three residues
$\overline\gamma(s),1,\overline\gamma(s)^{-1}$
are pairwise distinct.
The flag implies that the eigenvalues of $r(s)$
are $\gamma(s),1,\gamma(s)^{-1}$.
Reduction and pairwise distinctness force
\[
                a_0(s)=\gamma(s),\qquad
                a_1(s)=1,\qquad
                a_2(s)=\gamma(s)^{-1}.                       \tag{G55.7}
\]
There is no remaining permutation: each diagonal
entry has the residue of the corresponding entry
in (G55.6).

An invariant complete flag for this diagonal matrix
with these ordered, distinct eigenvalues has first
line $Le_0$ and first plane $Le_0\oplus Le_1$.
Indeed an invariant subspace is the sum of the
eigenlines belonging to its eigenvalues; the first
one and first two eigenvalues of the ordered flag
specify exactly those subspaces.
Applying the ordered flag to every inertia element
therefore gives, coordinate by coordinate,
$a_0=\gamma$, $a_1=1$ and $a_2=\gamma^{-1}$ on the
whole chosen inertia group.
Since $\chi=1$ there, all $\delta_i$ are trivial
there.

This conclusion holds above every place of $K$
dividing $\ell$.  If there are two such places,
their inertia groups are exchanged by $c$.
The equality $r(c)=r_\chi(c)$ makes conjugation by
$c$ reverse and invert both diagonal triples,
and equality on one inertia group implies equality
on the other.  Conjugation by an element of $H$
does not alter any diagonal character.
Thus all $\delta_i$ are everywhere unramified
characters of $G_K$.

\proofstep{The fixed class group kills the global ratios.}
An everywhere-unramified continuous abelian
character of $G_K$ factors through
$\operatorname{Cl}(\mathcal O_K)$, by class field
theory; $K$ has no real places.
Here $\mathcal O_K=\mathbb Z[\omega]$ has class
number one, as already used in AC-K04 and AC-G22.
One elementary verification is norm-Euclidean
division: after rounding the two coefficients
of a complex number in the basis $(1,\omega)$,
the difference is $u+v\omega$ with
$|u|,|v|\le1/2$, and
$u^2-uv+v^2\le3/4<1$.
This gives division with smaller norm and makes
every ideal principal.
Consequently each $\delta_i$ is trivial.
Equations (G55.3)--(G55.4) prove $r|_H=r_\chi|_H$.
Both representations have the same normalized
value $(I,-1)$ at $c$, so they agree on $\Gamma$.

Finally $f$ and $\overline{x}_0$ agree on every
actual matrix entry and on $O$.
They are continuous, by \texttt{Loc.ptCont} and
the finite-coefficient conclusion above, and
these entries generate an adically dense
$O$-algebra in $S$ by AC-G51.
Hence $f=\overline{x}_0$, proving both assertions
in (G55.2).
\end{proof}

\begin{corollary}[AC-G55-M3: the point condition conditional on the actual \texttt{M2}]
\label{AC-G55-M3}
For the actual model of AC-G51, \texttt{M2} implies
the complete source condition \texttt{M3}.
In fact the origin is the only point on
$Q.\mathrm{ideal}$ which annihilates every
nonzero weight part.
This result leaves \texttt{M2} and the mixed
two-sided point condition \texttt{M4} open.
\end{corollary}
\begin{proof}
Let $\xi:A\to\overline{\mathbb Q}_\ell$ vanish on
$Q.\mathrm{ideal}$ and on every positive and
negative weight part.
The actual \texttt{Model.pts} field, using
precisely the assumed \texttt{M2}, gives a torus
element $t$ and an actual local point $f$ of $S$
with
\[
                    t\cdot\xi=f\circ q_A.                    \tag{G55.8}
\]
The source field is stated this way in
\path{UnitaryThreeGroups/ConjStable.lean};
it is not an equality without a possible torus
translation at a general point.
Here all nonzero weight parts vanish at $\xi$,
so $t\cdot\xi=\xi$.
Every off-diagonal actual matrix entry has a
homogeneous preimage of nonzero weight in $A$.
Its value under $f$ is therefore zero.
This makes $\operatorname{push}(f,r^{\rm univ})$
both upper and lower, and AC-G55 implies
$f=\overline{x}_0$.
Equation (G55.8) gives $\xi=\overline{x}$.

In the literal \texttt{M3} definition in
\path{CommutativeAlgebra/OrbitVanishingLocal.lean},
the antecedent requires a point which vanishes
on all negative and positive weights but differs
from the origin.  No such point exists.
Thus \texttt{M3} holds.
This is the applicable fixed-point part of
\path{UnitaryThreeGroups/GradedModelPointConditionOfPtHyp.lean};
neither full \texttt{PtHyp} nor any assertion about
the tangent at another point was inserted.
\end{proof}

\paragraph{Scope and geometric dictionary.}
The new conclusion is uniqueness of the actual
diagonal specialization and the resulting
\texttt{M3} implication.  Together with AC-G53,
the fixed and basepoint-rigidity fields are
proved in the stated local range.
The remaining point-rigidity field, its mixed
global cohomology, and finite degree-zero
support are distinct requirements.
The finite unramified characters and covers
are interpreted using only the previously
audited SGA~1, Expos\'e~V, \S\S4 and~7 framework.
Their triviality here comes from the displayed
inertia calculation and class field theory,
not from a new geometric purity or motivic
existence assertion.
